% ===== 原文 pp.13–15（§2.5 问题集 I；§3 开头；§3.1 Hitchin 纤维化与 Teichmüller 分支） =====

% [原文 p.13]
\subsection{问题集 I}

本问题集意在使读者熟悉以各种不同方式编码 Higgs 丛定义数据的做法。其中的问题既有课堂上讲过的，也有一些额外的补充。

\begin{prob}{1}
\textbf{复 Higgs 丛。}如前所述，通过定义 2.9 可以定义主 $G_{\C}$-Higgs 丛，它们可以用一对 $(E,\Phi)$
来描述——其中 $E$ 是向量丛，$\Phi\colon E\to E\otimes K$ 是全纯 Higgs 场——并带有若干反映该群
性质的额外条件：
\begin{enumerate}[label=(\roman*),leftmargin=2.2em]
\item 给出用定义 2.1 中那样的对 $(E,\Phi)$ 描述 $\SO(2p+1,\C)$-Higgs 丛的方法，并列出需要满足的
  额外条件。
\item 证明：若 $\lambda$ 是 $\Sp(2p,\C)$-Higgs 丛 $(E,\Phi)$ 的一个特征值，则 $-\lambda$ 也必是
  它的特征值。
\end{enumerate}
\end{prob}
% [Tutor Note: p.13 原文 "a holomorphic Higgs filed Φ : E → E ⊗ K" 中 "filed" 为 "field" 的拼写错误
%  （同一段上文提及 Higgs bundle 的定义，且定义 2.1 中该对象即 Higgs 场）；中译按 "Higgs 场" 处理。]

\begin{prob}{2}
\textbf{实 Higgs 丛。}为了理解 $\Sp(2p,2q)$-Higgs 丛，我们首先考虑相应李代数的结构：
\begin{enumerate}[label=(\roman*),leftmargin=2.2em]
\item 给出复辛李代数的非紧实形式 $\fsp(2p,2q)$ 的一个 $4\times4$ 矩阵表达式（写出矩阵元素）。
\item 证明：对如例 2.14 中的 $\fsp(2p,2q)=\fh\oplus\fm$，可以把 $\fm_{\C}$ 表示为某些非对角矩阵
  构成的子集。
\item 证明：若 Higgs 丛 $(E,\Phi)$ 稳定，则对 $\lambda\in\C^{*}$ 与 $E$ 的全纯自同构 $\alpha$，
  诱导出的 Higgs 丛 $(E,\lambda\Phi)$ 与 $(E,\alpha^{*}\Phi)$ 也都稳定。
\item 把 $\Sp(2p,2q)$-Higgs 丛定义为满足各种反映该群性质的额外条件的对 $(E,\Phi)$。
\item 描述一般的 $\Sp(2p,2q)$-Higgs 场的特征多项式。
\end{enumerate}
\end{prob}
% [Tutor Note: p.13 问题 2(i) 原文说 "Give a 4 × 4 matrix expression ... of the non-compact real form
%  sp(2p, 2q)"，但 sp(2p,2q) ⊂ sp(2p+2q,C) 一般是 (2p+2q)×(2p+2q) 矩阵；"4 × 4" 只在 p=q=1 时成立
%  （sp(2,2) ⊂ sp(4,C)）。正文照原文排，读者可先取 p=q=1 作具体计算。]

\begin{prob}{3}
\textbf{实 Higgs 丛。}设 $\fg_{\C}$ 是带有复结构 $\ii$ 的复李代数，其李群为 $G_{\C}$。回忆：$\fg_{\C}$
的一个\emph{实形式}是满足 $\fg_{\C}=\fg\oplus\ii\fg$ 的实李代数。给定 $\fg_{\C}$ 的实形式 $\fg$，
元素 $Z\in\fg_{\C}$ 可以写成 $Z=X+\ii Y$，其中 $X,Y\in\fg$。映射 $X+\ii Y\mapsto X-\ii Y$ 称为
关于 $\fg$ 的\emph{共轭}。此外，$\fg_{\C}$ 的任一实形式都由 $\fg_{\C}$ 上某个反线性对合 $\tau$ 的
不动点集给出，即由满足
\[
\begin{aligned}
\tau(\tau(X))&=X, & \tau(zX)&=\bar z\,\tau(X),\\
\tau(X+Y)&=\tau(X)+\tau(Y), & \tau([X,Y])&=[\tau(X),\tau(Y)]
\end{aligned}
\]
的映射给出，其中 $X,Y\in\fg_{\C}$，$z\in\C$。特别地，关于 $\fg$ 的共轭满足这些性质。由上述定义，
完成以下各题：
\begin{enumerate}[label=(\roman*),leftmargin=2.2em]
\item 求出 $\fg_{\C}$ 上的一个对合 $\sigma$，使得：
  \begin{itemize}[label=--,leftmargin=2.6em]
  \item $\sigma$ 的不动点集 $\fg^{\sigma}_{\C}$ 由 $\fg$ 的极大紧子代数 $\fh$ 的复化给出；
  \item 该对合 $\sigma$ 下的反不变集由 $\fm_{\C}$ 给出，其中 $\fg=\fh\oplus\fm$ 是 Cartan 分解。
  \end{itemize}
% [原文 p.14]
\item 给出 $\mathfrak{gl}(n,\C)$ 与 $\fsl(n,\C)$ 上固定号差为 $(p,q)$ 的酉群 $\fu(p,q)$ 与
  $\fsu(p,q)$ 的对合的描述。
\item 把 $\Sp(2p,2q)$-Higgs 丛描述为整个 $\Sp(2p+2q,\C)$-Higgs 丛模空间上某个对合的不动点集：该对
  合由李代数上的作用诱导，并用到 Hitchin 方程与上面的 \textbf{问题 2}。
\end{enumerate}
\end{prob}
% [Tutor Note: 这里的四个条件是反线性对合 τ 的定义性质，缺一不可：τ∘τ = id（对合）、反线性
%  τ(zX) = z̄ τ(X)、实线性 τ(X+Y) = τ(X)+τ(Y)、以及保持李括号（故为李代数自同构）。两条实线性条件
%  合起来等价于 τ 是 R-线性的，配上反线性条件即"反线性";固定点集 {X : τ(X)=X} 是实子空间，
%  其复化回到 g_C。]

\begin{prob}{4}
\textbf{抛物 Higgs 丛。}考虑 $\Gl(2,\C)$-Higgs 丛时，回忆 Hitchin 在 [61] 中建立了 $\pi_1(\Sigma)$
的不可约 $\Gl(2,\C)$ 表示的等价类与次数为零的 $\Gl(2,\C)$-Higgs 丛的同构类之间的一个对应。此后
Simpson 在 [33] 中得到了抛物 Higgs 丛与 $\Sigma/\{p_1,\dots,p_n\}$ 的基本群之间类似的对应。通过
这些对应，可以获得抛物 Higgs 丛及其对应表示的某些特殊性质，正如例如 [21] 中所做的那样。
\begin{enumerate}[label=(\roman*),leftmargin=2.2em]
\item 给出一个次数为零的秩 2 抛物稳定 Higgs 丛 $(E,\Phi)$ 的例子。
\item 利用例 2.16 中的描述，给出一个抛物稳定 $\U(3,1)$-Higgs 丛 $(E,\Phi)$ 的例子。
\item 给出 \textbf{问题 4.i.} 中 Higgs 丛 $(E,\Phi)$ 所对应的平坦联络的表达式，并修改该对，直到能够
  证明该平坦联络的和乐（经共轭后）包含在 $\Sl(2,\R)$ 中。
\item 由 [42] 回忆：$E$ 上的一个\emph{最小抛物结构}是指在每个 $p_i$ 处选定一个部分旗
  $0\subset L_i\subset E_{p_i}$，其中 $L_i$ 是一条直线。一个\emph{最小抛物 Higgs 场}是一个
  $\cO_X$-线性映射 $\phi\in H^0(X,\End(E)\otimes K_X(D))$，使得在每个点 $p_i$ 处，$\phi$ 的残数
  $\phi_i$ 关于该旗是严格三角的：$\phi(E_{p_i})\subseteq L_i$，$\phi_i(L)=0$。构造一个秩 3 的
  \emph{最小抛物 Higgs 场}的例子。
\item 证明：若 \textbf{问题 4.iv.} 中的最小抛物 Higgs 场 $\phi$ 是严格抛物的，则残数 $\phi_i$
  是 2 阶幂零的，并且落在 $\fsl_r$ 的最小幂零轨道的闭包中。
\end{enumerate}
\end{prob}
% [Tutor Note: p.14 问题 4 开篇有两处原文笔误：① "parabolic Higgs bundles and and the fundamental
%  group" 中 "and" 重复；② "Though those identifications, ... can be obtained" 应为 "Through those
%  identifications"，否则句子无法解析（中文按 "通过这些对应" 译出）。正文不改动原文，仅在此说明。]

\begin{prob}{5}
\textbf{野 Higgs 丛。}考虑一个秩 2 的 Higgs 丛，它在 $p_2$ 处有一个 3 阶极点、在 $p_1$ 处有一个
1 阶极点。为此，取对角矩阵 ${}^{j}A_i$，以及
\[
\nabla := \dd - \left({}^{1}A_1\frac{\dd z}{(z-p_1)}\right)
 - \left({}^{2}A_3\frac{\dd z}{(z-p_2)^3} + {}^{2}A_2\frac{\dd z}{(z-p_2)^2}
 + {}^{2}A_1\frac{\dd z}{(z-p_2)}\right).
\]
利用本章中给出的例子，考虑以下问题：
\begin{enumerate}[label=(\roman*),leftmargin=2.2em]
\item 描述反 Stokes 方向的集合，并判断 Stokes 矩阵是上三角的还是下三角的；
\item 如图 2.24 那样给出触手的描述。
\item 对如注 2.22 中那样带有一个 4 阶极点的 Higgs 丛，证明该系统有 6 个反 Stokes 方向，并给出
  它们的描述。
\end{enumerate}
\end{prob}
% [Tutor Note: 记号 ${}^{j}A_i$ 的读法：左上标 j 标记极点 $p_j$（此处 j = 1,2），下标 i 标记该极点
%  处展开的阶数。故第一个括号是 p_1 处的 1 阶极点项，其系数为 ${}^{1}A_1$；第二个括号是 p_2 处的
%  3 阶极点项，系数从最高阶到最低阶依次为 ${}^{2}A_3,{}^{2}A_2,{}^{2}A_1$——与题干 "order 3 at p_2,
%  order 1 at p_1" 一致，可用来核对阶数与上标是否对应。]

% [原文 p.15]
\section{Hitchin 纤维化的几何}

\begin{quote}
\kaishu
一个微分方程组的可积性，应当通过某些一般可辨认的特征表现出来：
\begin{itemize}[label=$\star$,leftmargin=1.6em,itemsep=1pt,topsep=2pt]
\item 存在大量守恒量；
\item 存在代数几何；
\item 能够给出显式解。
\end{itemize}
这些准则应在非常宽泛的意义上理解。”
\end{quote}
\begin{flushright}
——Nigel Hitchin
\end{flushright}
% [Tutor Note: 原文这段题词（p.15）只有结尾的一个右引号 ”（"a very broad sense.\""），并无起首引号；
%  本译本照原样保留这一单独的右引号，不代作者补全。]

$G_{\C}$-Higgs 丛的模空间 $\mathcal{M}_{G_{\C}}$ 可以自然地通过 Hitchin 纤维化来研究，该纤维化由
[61] 引入。本章将专门研究此纤维化，并描述它的一些几何性质。

\subsection{Hitchin 纤维化与 Teichmüller 分支}

该 Hitchin 纤维化是通过 $\fg_{\C}$（即 $G_{\C}$ 的李代数）的不变多项式代数的一组齐次基
$p_i$（$i=1,\dots,k$）来定义的，这些基元的次数记作 $d_i$。于是，Hitchin 纤维化由
\[
\begin{gathered}
h:\mathcal{M}_{G_{\C}}\longrightarrow\mathcal{A}_{G_{\C}}
 :=\bigoplus_{i=1}^{k}H^0(\Sigma,K^{d_i}),\\
(E,\Phi)\mapsto\big(p_1(\Phi),\dots,p_k(\Phi)\big).
\end{gathered}
\]
给出。映射 $h$ 称为 Hitchin 映射，它对任意一组基的选择都是紧映射（proper map）（细节见
[61, Section 4]）。此外，Hitchin 底 $\mathcal{A}_{G_{\C}}$ 总满足
\[
\dim\mathcal{A}_{G_{\C}}=\dim\mathcal{M}_{G_{\C}}/2,
\]
从而使 Higgs 丛模空间成为一个可积系统。

通过选取不变多项式的不同齐次基，可以使 Hitchin 纤维化几何的不同侧面显现出来。特别地，人们有时
会问：纤维化的某个性质（或某条纤维的某个性质）在基变换之下是否保持不变；我们将在后面几节看到
这样的例子。在整份讲义中，对于秩为 $n$ 的经典 Higgs 丛 $(E,\Phi)$，我们经常取由 $\tr(\Phi^i)$
（$1\leqslant i\leqslant n$）给出的基，或者取由 Higgs 场 $\Phi$ 的特征多项式的系数给出的基。

\begin{rem}{3.1}
通过基变换，可以得到同一个群的 Hitchin 纤维化的不同表达式之间的对应。例如，特征多项式的系数
可以用迹表示如下：
\begin{equation}
p(x)=x^n+\tr(\Phi)x^{n-1}+\sum_{k=2}^{n-1}x^{n-k}(-1)^k\Tr(\Lambda^k\Phi)+(-1)^n\det(\Phi) \tag{3.2}
\end{equation}
其中 $\Tr(\Lambda^k\Phi)$ 是 $\Phi$ 的 $k$ 次外幂的迹，其维数为 $\dfrac{n!}{k!\,(n-k)!}$
\end{rem}
% [Tutor Note: p.15 式 (3.2) 首项符号有误：按通项模式 (−1)^k 与 p.41 同一展开式，k=1 项应为
%  $-\tr(\Phi)x^{n-1}$，原文印作 $+\tr(\Phi)x^{n-1}$（正文照原文排）。正确形式：
%  $p(x)=x^n-\tr(\Phi)x^{n-1}+\sum_{k=2}^{n-1}(-1)^k x^{n-k}\Tr(\Lambda^k\Phi)+(-1)^n\det(\Phi)$。
%  验算：$n=2$ 时 $p(x)=x^2-\tr(\Phi)x+\det(\Phi)$；$n=3$ 时 $p(x)=x^3-\tr(\Phi)x^2+\Tr(\Lambda^2\Phi)x-\det(\Phi)$。
%  另：原文首项写 tr、通项写 Tr，大小写不统一，中译按原样分别用 \tr 与 \Tr。]
% [Tutor Note: p.15 末 Remark 3.1 以 "with dimension n!/(k!(n−k)!)" 结束，原文该分式后没有句号
%  （p.16 另起新段）；中译按句末处理，补出句号。]
